\chapter{Computing a BNE is at least as hard as ARRIVAL}\label{ch:arrival}

\paragraph{Idea.} By Chapter~\ref{ch:market}, a market with slopes at least one can be hidden in a
one-resource game. So it suffices to build, from an ARRIVAL instance, a market whose every equilibrium
reveals where the run ends. The bridge is the continuous switch map of Section~\ref{sec:pre-arrival}:
each switch vertex $v$ must split the flow $x_v$ entering it into $h_0(x_v)$ and $h_1(x_v)$. We build a
market \emph{gadget} that does exactly this. Inside the gadget, min-cost flow does the work: two copies
of the Disser--Skutella counting network share a source, and the cheapest way to send $x$ units
alternates between the copies, unit by unit, like a switch.

The project's AI agents proposed the one-resource reduction and designed
the switch using two copies of the counting network. In a later comment
on the manuscript, quoted here from Lin's record, Conitzer wrote
``Incidentally someone else suggested ARRIVAL here'' and linked to
Jeonseong Kim's post~\cite{kim26}.
Theorem~\ref{thm:A} specifies the transformation and proves correctness
for every equilibrium of its output game.

\begin{theorem}[ARRIVAL-hardness]\label{thm:A}
\statusrefereed
There are polynomial-time algorithms $\Red$ and $\Dec$ with the following properties. $\Red$ maps an
ARRIVAL instance with $n$ switch vertices to an explicit Bayesian simple security game $\Gamma$ with one
defender resource, at most $12n^2+6n+6$ attacker types and at most $12n^2+6n+7$ targets, in which the
defender gets $0$ at every uncovered target and $1$ at every covered target. For every BNE $\beta$ of
$\Gamma$, $\Dec(\beta)$ is the correct answer to the ARRIVAL instance.
\end{theorem}

\begin{corollary}\label{cor:A}
\statusrefereed
If a BNE of such games can be computed in polynomial time, then ARRIVAL is in P.
\end{corollary}

The same hardness result, for the shorter construction of Appendix~\ref{app:short}, is machine-checked in Lean for rational BNEs, except that $\Red$ and $\Dec$ run in polynomial time (Section~\ref{sec:lean-status}); the construction of this chapter is not formalized.


Theorem~\ref{thm:A} rules out a polynomial-time BNE algorithm only if ARRIVAL is not in $\mathsf P$, and
it says nothing against subexponential algorithms. It is much weaker than $\mathsf{PPAD}$-hardness and does not imply it. It
concerns \emph{exact} BNE: the decoder uses the exact equilibrium conditions. Since the search version of ARRIVAL is in
$\mathsf{UEOPL}$, every $\mathsf{UEOPL}$-hard problem is ARRIVAL-hard, and natural $\mathsf{UEOPL}$-hard problems are
known~\cite{hmr26}, so this is not the first natural ARRIVAL-hard problem; it is a direct reduction
from ARRIVAL to a problem that is not known to be $\mathsf{UEOPL}$-hard.

Appendix~\ref{app:short} gives a shorter proof of Theorem~\ref{thm:A} by a different construction. This chapter's construction is the one Chapter~\ref{ch:upper} builds on: its market (Theorem~\ref{thm:arrival-m1}) gives ARRIVAL $\le_m$ P-LCP (Corollary~\ref{plcp:cor:arrival}).

\paragraph{Structure of the proof.} The proof has three links.
\begin{enumerate}[nosep]
\item An (M1) market whose slopes are all at least $1$ embeds into a one-resource game, and every BNE
decodes to a market equilibrium. This is Theorem~\ref{thm:mk-embed} with $\omega=1$.
\item An ARRIVAL instance is composed into one (M1) market from copies of a switch gadget. If the gadget
meets a local specification (S1--S4), every market equilibrium is a fixed point of the switch map, and
every such fixed point puts load $1$ on the terminal the run reaches (Section~\ref{sec:ar-compose}).
\item A gadget with $12\log_2K-6$ buyers and goods and $O(\log^2K)$ bits meets S1 and S2 for $K=2^{n+1}$ (S3 with the composition's choice of $c_{p_v}$; S4 by construction)
(Section~\ref{sec:ar-gadget}).
\end{enumerate}
The decoder $\Dec$ maps the BNE to a market equilibrium and reads the load on the \yes{} terminal good. Figure~\ref{fig:ar-map} maps the proof.

\begin{figure}[p]
\centering
\begin{tikzpicture}[
  font=\footnotesize,
  nohy/.style={execute at begin node={\hyphenpenalty=10000\exhyphenpenalty=10000}},
  res/.style={draw,rounded corners,align=left,inner sep=3pt,nohy},
  cit/.style={draw,dashed,gray!80!black,text=black!80,rounded corners,align=left,inner sep=3pt,nohy},
  obj/.style={draw,rounded corners,align=center,inner sep=3pt,fill=gray!8,text width=3.0cm,
    minimum height=1.0cm,nohy},
  use/.style={-{Stealth[length=2mm]},thick},
  usec/.style={-{Stealth[length=2mm]},thick,gray!80!black},
  lb/.style={font=\scriptsize,fill=white,inner sep=1pt,align=center,nohy}]
% ---------------- (a) the reduction and the decoder ----------------
\node[anchor=west,font=\small\bfseries] at (0,1.2) {(a) What $\Red$ builds, and how $\Dec$ reads a BNE};
\node[obj] (inst) at (1.6,0)    {ARRIVAL instance with $n$ switch vertices};
\node[obj] (mkt)  at (8.0,0)    {(M1) market with $O(n^2)$ buyers and goods};
\node[obj] (game) at (14.2,0)   {game $\Gamma_1$ with one resource, $O(n^2)$ types and targets};
\node[obj] (ans)  at (1.6,-2.4) {answer \yes{} if and only if that load is $1$};
\node[obj] (meq)  at (8.0,-2.4) {an equilibrium of the market};
\node[obj] (bne)  at (14.2,-2.4){any BNE $\beta$ of $\Gamma_1$};
\draw[use] (inst) -- node[lb,above=1pt,text width=2.8cm]{$\Red$: preprocess, then one switch gadget per vertex}
  node[lb,below=1pt]{Theorem~\ref{thm:arrival-m1}} (mkt);
\draw[use] (mkt) -- node[lb,above=1pt,text width=2.6cm]{$\Red$: embed with $\omega=1$}
  node[lb,below=1pt]{Theorem~\ref{thm:mk-embed}} (game);
\draw[use,dashed] (game) -- node[lb,right=2pt,fill=none,align=left]{any BNE\\algorithm} (bne);
\draw[use] (bne) -- node[lb,above=1pt,text width=2.6cm]{$\Dec$: read prices and purchases off $\beta$}
  node[lb,below=1pt]{Theorem~\ref{thm:mk-embed}} (meq);
\draw[use] (meq) -- node[lb,above=1pt,text width=2.8cm]{$\Dec$: load of non-terminal buyers on $T_{\yes}$}
  node[lb,below=1pt]{correct by Theorem~\ref{thm:arrival-m1}} (ans);
% ---------------- (b) the dependencies ----------------
\node[anchor=west,font=\small\bfseries] at (0,-3.95) {(b) Which result is used where};
\node[res,very thick,text width=7.2cm,align=center] (A) at (7.9,-4.85)
  {\textbf{Theorem~\ref{thm:A}} (ARRIVAL-hardness): $\Red$ and $\Dec$ run in polynomial time, and
  $\Dec(\beta)$ is the correct answer for every BNE $\beta$};
\node[res,text width=5.5cm,anchor=north] (E) at (2.85,-6.0)
  {\textbf{Theorem~\ref{thm:mk-embed}} (Chapter~\ref{ch:market}): a market whose slopes are all $\ge1$
  embeds in a one-resource game, and every BNE of that game gives a market equilibrium};
\node[res,text width=7.2cm,anchor=north] (AM) at (12.1,-6.0)
  {\textbf{Theorem~\ref{thm:arrival-m1}} (ARRIVAL to (M1)): in every equilibrium of the composed market,
  the non-terminal buyers put load $1$ on the terminal that the run reaches and $0$ on the other.
  S3 and S4 hold by construction.};
\node[res,text width=7.0cm,anchor=north] (C) at ($(3.6,0 |- E.south)+(0,-0.9)$)
  {\textbf{Theorem~\ref{thm:compose}} (Composition): if the gadget meets S1--S4, then in every
  equilibrium the port prices lie in $[0,\bar P]$, halving gives $x_v\le K$, the loads form a fixed
  point, and the terminal that the run reaches gets load $1$};
\node[res,text width=7.2cm,anchor=north] (G) at ($(AM.south)+(0,-0.9)$)
  {\textbf{Theorem~\ref{thm:gadget}} (the gadget meets S1 and S2), from C1--C6:\\[1pt]
  $\bullet$ Lemma~\ref{lem:lp}: at fixed port prices the gadget market is an LP\\
  $\bullet$ Lemma~\ref{lem:full} (C3): every private good is full\\
  $\bullet$ Lemma~\ref{lem:flow} (C3--C6): every optimum is a flow of value $\min(x,K)$\\
  $\bullet$ Lemma~\ref{lem:alt} (C2, so $0<\theta<6$): min-cost flows send $h_0(x)$ to $a$ and
  $h_1(x)$ to $b$\\
  $\bullet$ C1 gives the ratio bound in S1};
\node[res,text width=7.0cm,anchor=north] (D) at ($(C.south)+(0,-0.8)$)
  {\textbf{Lemma~\ref{lem:decoder}} (Decoder): every fixed point puts inflow $1$ on the terminal that
  the run reaches and $0$ on the other};
\node[cit,text width=7.0cm,anchor=north] (H) at ($(D.south)+(0,-0.7)$)
  {cited: Haslebacher~\cite{haslebacher25}, Lemma~14 and Theorem~1; Dohrau et al.~\cite{dgkmw17},
  Lemma~1 (switching flows certify where the run ends)};
\node[res,text width=7.2cm,anchor=north] (I) at ($(G.south)+(0,-0.9)$)
  {\textbf{Lemma~\ref{lem:incr}} (Increments): the unit costs of the counting network $G_k$ have gaps
  $\ge6$, and its maximum flow is $2^k$};
\node[cit,text width=7.2cm,anchor=north] (DS) at ($(I.south)+(0,-0.7)$)
  {cited: Disser and Skutella~\cite{ds15}, the counting network (\S3.1) and the alternation argument
  that Lemma~\ref{lem:alt} adapts (\S3.2)};
% edges: an arrow goes from a result to a proof that uses it
\draw[use] (E.north) |- node[lb,pos=0.3,text width=1.9cm]{$\omega=1$; slopes $1$ and $\lambda$} (A.west);
\draw[use] (10.3,0 |- AM.north) -- (10.3,0 |- A.south);
\draw[use] (6.6,0 |- C.north) |- (AM.west);
\draw[use] (G.north) -- node[lb]{S1, S2} (AM.south);
\draw[use] (D.north) -- node[lb]{step (d)} (C.south);
\draw[usec] (H.north) -- (D.south);
\draw[use] (I.north) -- node[lb]{gaps: Lemma~\ref{lem:alt}; maximum flow: Lemma~\ref{lem:flow}} (G.south);
\draw[usec] (DS.north) -- node[lb]{the network $G_k$} (I.south);
\draw[usec] (DS.west) -- ++(-0.4,0) |- node[pos=0.3,sloped,above,inner sep=1pt,font=\scriptsize,text=gray]{alternation (Lemma~\ref{lem:alt})} (G.west);
\end{tikzpicture}
\caption{A map of the proof of Theorem~\ref{thm:A}. (a) What $\Red$ builds, and the two steps of $\Dec$;
exact sizes are in Section~\ref{sec:ar-size}. (b) The results that the proof uses. An arrow points from
a result to a proof that uses it, and dashed gray boxes are cited results. The local specification
S1--S4 (Section~\ref{sec:ar-compose}) connects the left and right halves: Theorem~\ref{thm:compose}
assumes it; Theorem~\ref{thm:gadget} supplies S1 and S2, and the construction in
Section~\ref{sec:ar-gadget-market} supplies S3 and S4. Conditions C1--C6 (also in
Section~\ref{sec:ar-gadget-market}) are inequalities on the gadget's parameters, and the labels show where the main ones are used. From Chapter~\ref{ch:prelim} the proof uses the game (Definition~\ref{def:game}),
the market (M1) (Definition~\ref{def:m1}) and Section~\ref{sec:pre-arrival}: ARRIVAL, its
preprocessing, the switch map $h_0,h_1$ and the fixed-point equations~\eqref{eq:fixedpoint}. The chapter proves the results
bottom-up: Lemma~\ref{lem:decoder}, Theorem~\ref{thm:compose}, the gadget
(Section~\ref{sec:ar-gadget}), then Theorem~\ref{thm:arrival-m1} and Theorem~\ref{thm:A}.}
\label{fig:ar-map}
\end{figure}


\section{Fixed points of the switch map decide ARRIVAL}

\begin{lemma}[Decoder]\label{lem:decoder}
\statusrefereed
In every fixed point~\eqref{eq:fixedpoint}, the terminal that the run reaches has inflow $1$ and the
other terminal has inflow $0$.
\end{lemma}

\begin{proof}
The proof has three steps: make the fixed point integral, observe that an integral fixed point is a
switching flow, and invoke the known theory of switching flows.

\emph{Step 1 (removing fractional vertices).} If $x_u$ is an integer, both out-edges of $u$ are
integral. On $(2z,2z+1)$ only the label-$0$ edge is fractional, and on $(2z+1,2z+2)$ only the label-$1$
edge, with slope $1$ on that piece. A fractional $x_u$ needs a fractional in-edge, because the source
term is integral, and the tail of that in-edge is fractional. Walking backwards therefore closes a simple
cycle $C$ of fractional vertices, and each edge of $C$ is the unique fractional out-edge of its tail.
Terminals have no out-edges, so they lie on no such cycle. Lower $x$ on $C$ by $\Delta$, the least
fractional part on $C$. Each vertex of $C$ stays on its closed affine piece. Its active out-edge, which is
its $C$-edge, drops by $\Delta$, and so does its $C$-in-edge. Its inactive edge is constant, and vertices
off $C$ are unchanged. The result is again a fixed point with the same terminal inflows and at least one
fewer fractional vertex. After at most $n$ rounds the fixed point is integral. (The backward walk is the
first part of~\cite[Lemma 14]{haslebacher25}; the lowering step is new.)

\emph{Step 2 (integral fixed points).} An integral fixed point is a switching
flow~\cite[Lemma 14]{haslebacher25}; this is also immediate from the values of $h_0$ and $h_1$ at
integers. Its terminal inflows are integers summing to $1$, so they are $(1,0)$ or $(0,1)$.

\emph{Step 3 (the run).} Suppose the unit sits on terminal $t$, and let $t'$ be the other terminal. Give
both $t$ and $t'$ two self-loops carrying flow $0$ (in~\cite{dgkmw17} the successor maps are defined at every vertex, the destination included, and the switching-flow conditions must hold there too; with these loops they hold at $t$ and $t'$ trivially). The fixed point is then a switching flow for destination $t$ in
the sense of~\cite{dgkmw17}. By~\cite[Lemma 1]{dgkmw17}, if a switching flow exists then
$\mathrm{Run}(G,o,t)$ terminates. The run therefore reaches $t$, and it reaches $t$ before $t'$, since a
token at $t'$ stays there. (Alternatively, \cite[Theorem 1]{haslebacher25} states the needed fact
directly.)
\end{proof}

\section{Gadgets and composition}\label{sec:ar-compose}

\paragraph{Idea.} A gadget is a small market with one input good $p$ and two output ``ports'' $a$ and
$b$. Its specification says: if $x$ units of outside demand sit on $p$, the gadget sends exactly
$h_0(x)$ to port $a$ and $h_1(x)$ to port $b$, whatever the port prices are within a box. Gluing a gadget
per vertex, with ports equal to the input goods of the successors, turns market equilibria into fixed
points of~\eqref{eq:fixedpoint}.

\subsection{The local specification}

A \emph{gadget} consists of an input good $p$ with capacity $c_p$, private goods, internal buyers, and two
\emph{port goods} $a$ and $b$ that lie outside it. An internal option is \emph{$p$-labelled},
\emph{$a$-labelled} or \emph{$b$-labelled} according to which of these goods it targets. Labels matter
because $p$, $a$ and $b$ may coincide after composition. A \emph{gadget equilibrium at $(x,P_a,P_b)$} is
defined for an external load $x\ge0$ on $p$ and given port prices $P_a,P_b\ge0$. It requires conditions
(a)--(d) of Definition~\ref{def:m1} for the internal buyers and private goods, and (a) and (b) for port
options at the given prices. For $p$ it requires ($p$-labelled internal load)${}+x\le c_p$, with
equality if $P_p>0$. The ports have no capacity condition. Write $A$ and $B$ for the internal
$a$-labelled and $b$-labelled loads. Fix $K=2^{n+1}$ and a price bound $\bar P>0$.

\begin{description}[leftmargin=2.5em,style=nextline]
\item[S1.] For all $x\in[0,K]$ and $(P_a,P_b)\in[0,\bar P]^2$, every gadget equilibrium has $A=h_0(x)$
and $B=h_1(x)$. Every $p$-labelled option $(i,p)$ has $a_{ip}/b_{ip}\le\bar P$.
\item[S2.] For all $x\in[K,c_p)$ and $(P_a,P_b)\in[0,\bar P]^2$, every gadget equilibrium has
$A,B\ge K/2$ and $A+B\le x$.
\item[S3.] In the composed market, $c_p$ exceeds the total mass of all buyers that hold a
non-$p$-labelled option on this gadget's $p$. These are the predecessors' port buyers, this gadget's own
port buyers when it has a self-loop, and the dispatch buyer.
\item[S4.] No internal buyer holds options on two of $p$, $a$ and $b$.
\end{description}
Without S4, a self-loop would give one buyer two options on the same good, which (M1) does not allow.

\subsection{Composition}

Given a preprocessed instance (Section~\ref{sec:pre-arrival}) and a gadget, the composed market
consists of:
\begin{itemize}[nosep]
\item one copy of the gadget for each $v\in V$, with input good $p_v$. Its $a$-port is the input good of
$\succz(v)$, or the terminal good if $\succz(v)$ is a terminal. Its $b$-port is defined in the same way
from $\succo(v)$;
\item a \emph{dispatch} buyer of mass $1$ whose only option, on $p_o$, has intercept-to-slope ratio above $\bar P$;
\item for each terminal $t$, a good $T_t$ whose capacity $C_T$ exceeds the total mass of the
non-terminal buyers that hold an option on $T_t$, and a \emph{terminal buyer} of mass $>C_T$ whose only
option is $T_t$, with ratio $r_T\in(0,\bar P)$.
\end{itemize}

Figure~\ref{fig:ar-instance} shows the composition for a two-vertex instance.

\begin{figure}[t]
\centering
\begin{tikzpicture}[
  sw/.style={circle,draw,minimum size=7mm,font=\small},
  te/.style={draw,rounded corners,minimum size=6mm,font=\small},
  gd/.style={draw,thick,rounded corners,minimum width=2.0cm,minimum height=1.0cm,align=center,font=\small,fill=gray!8},
  good/.style={draw,minimum width=1.0cm,minimum height=0.6cm,font=\small},
  buy/.style={draw,rounded corners=6pt,font=\scriptsize,align=center,inner sep=2pt},
  ar/.style={-{Stealth[length=2mm]},thick},
  lb/.style={font=\scriptsize,fill=white,inner sep=1pt}]
  % switch graph
  \node[font=\small] at (-1.5,5.3) {(i)};
  \node[sw] (u) at (0,4.6) {$u$};
  \node[sw] (v) at (3,4.6) {$v$};
  \node[te] (Y) at (6,5.3) {\yes};
  \node[te] (N) at (6,3.9) {\no};
  \draw[ar] (-0.9,4.6) node[left,font=\scriptsize]{start} -- (u);
  \draw[ar] (u) -- node[lb,above]{$0$} (v);
  \draw[ar] (u) to[out=60,in=120,looseness=6] node[lb,above]{$1$} (u);
  \draw[ar] (v) -- node[lb,above]{$0$} (Y);
  \draw[ar] (v) -- node[lb,below]{$1$} (N);
  % market
  \node[font=\small] at (-2.3,2.2) {(ii)};
  \node[buy] (disp) at (-2.3,1.2) {dispatch\\mass $1$};
  \node[good] (pu) at (-0.4,1.2) {$p_u$};
  \node[gd] (Gu) at (2.2,1.2) {gadget of $u$};
  \node[good] (pv) at (5.6,1.2) {$p_v$};
  \node[gd] (Gv) at (8.0,1.2) {gadget of $v$};
  \node[good] (Ty) at (10.9,2.1) {$T_{\yes}$};
  \node[good] (Tn) at (10.9,0.3) {$T_{\no}$};
  \node[buy] (by) at (12.5,2.1) {terminal\\buyer};
  \node[buy] (bn) at (12.5,0.3) {terminal\\buyer};
  \draw[ar] (disp) -- node[lb,above]{$1$} (pu);
  \draw[ar] (pu) -- node[lb,above]{$x_u$} (Gu);
  \draw[ar] (Gu) -- node[lb,above]{$a$: $h_0(x_u)$} (pv);
  \draw[ar] (Gu.south) .. controls (2.2,-0.4) and (-0.4,-0.4) .. node[lb,below]{$b$: $h_1(x_u)$} (pu.south);
  \draw[ar] (pv) -- node[lb,above]{$x_v$} (Gv);
  \draw[ar] (Gv) -- node[lb,above,sloped,pos=0.45]{$a$: $1$} (Ty);
  \draw[ar] (Gv) -- node[lb,below,sloped,pos=0.45]{$b$: $0$} (Tn);
  \draw[thick] (by) -- (Ty);
  \draw[thick] (bn) -- (Tn);
\end{tikzpicture}
\caption{(i) An ARRIVAL instance with origin $u$: $\succz(u)=v$, $\succo(u)=u$, $\succz(v)=\yes$,
$\succo(v)=\no$. The run is $u\to v\to\yes$: the token leaves $u$ by exit~$0$ (and $u$'s switch turns to
$1$), then leaves $v$ by exit~$0$. The instance is discriminating: if both switches started at
$\succo$, the run would be $u\to u\to v\to\no$. (ii) The composed market. The gadget of $u$ has input
good $p_u$, $a$-port $p_v$ and, because of the self-loop, $b$-port $p_u$; the gadget of $v$ has
$a$-port $T_{\yes}$ and $b$-port $T_{\no}$. Labels show the loads in an equilibrium. By Theorem~\ref{thm:compose} they form a fixed point
of~\eqref{eq:fixedpoint}: $x_u=1+h_1(x_u)$ and $x_v=h_0(x_u)$. The first equation holds exactly for
$x_u\in[1,2]$ (on that interval $h_1(x_u)=x_u-1$), so fixed points need not be integral, but every one
has $x_v=h_0(x_u)=1$; hence $T_{\yes}$ receives $h_0(1)=1$ and $T_{\no}$ receives $h_1(1)=0$, the
terminal the run reaches (Lemma~\ref{lem:decoder}).}
\label{fig:ar-instance}
\end{figure}

\begin{theorem}[Composition]\label{thm:compose}
\statusrefereed
Assume S1--S4 and $K=2^{n+1}$. Every equilibrium of the composed market has the following properties:
\begin{enumerate}[nosep]
\item every port price lies in $[0,\bar P]$;
\item $x_v\le K$ for all $v$, where $x_v$ is the non-$p$-labelled load on $p_v$, and $x$ is a fixed
point;
\item the load that non-terminal buyers place on $T_t$ is $1$ at the terminal the run reaches and $0$ at
the other.
\end{enumerate}
\end{theorem}

\begin{proof}
The proof restricts the equilibrium to each gadget, bounds the prices, and then shows by a halving
argument that no vertex can receive more than $K$ units.

\emph{(a) Restriction.} Fix gadget $v$, and let $x_v$ be the load on $p_v$ through all options that are
not $p$-labelled options of $v$. Take the port prices from the market. Internal buyers keep their option
sets, private goods are bought only by internal buyers, and the condition on $p_v$ reads ($p$-labelled
internal load)${}+x_v\le c_p$, with equality if $P_{p_v}>0$. So the restriction is a gadget equilibrium
at $(x_v,P_a,P_b)$. This includes coincident ports. If $a=b=w$, then $P_a=P_b$, which is a point of the
box, and both labelled loads go to $p_w$. If $a=p_v$, the $a$-labelled load is part of $x_v$, and the
restriction is still a gadget equilibrium.

\emph{(b) Prices.} By S3, $x_v<c_p$. If $P_{p_v}>0$, then $p_v$ is full, so its $p$-labelled internal
load $c_p-x_v$ is positive. Some $p$-labelled buyer $i$ is then tight, which gives
$0\le v_i=a_{ip}-b_{ip}P_{p_v}$, so $P_{p_v}\le a_{ip}/b_{ip}\le\bar P$. At a terminal, $P_T<r_T$ would
give the terminal buyer positive value, and it would then demand more than $C_T$. $P_T>r_T$ would make
the terminal buyer leave, so $T$ would have to be filled by the other buyers, whose mass is below $C_T$.
So $P_T=r_T$. All port prices therefore lie in $[0,\bar P]$, and the dispatch buyer, with ratio
$>\bar P$, has positive value and buys exactly $1$.

\emph{(c) Load bound.} By S1 or S2 at each gadget, $A_v+B_v\le x_v$. Counting labelled loads gives total
terminal inflow ${}=1+\sum_v(A_v+B_v-x_v)\le1$. Every output satisfies output
${}\ge\phi(x)=\min((x-1)/2,K/2)$: where $x\le K$, the output is at least $h_1(x)\ge(x-1)/2$ (check this
on $[2z,2z+1]$ and on $[2z+1,2z+2]$); where $K\le x<c_p$, S2 gives at least $K/2$. Suppose $x_v>K$. Take
a shortest path $v=u_0,u_1,\dots,u_l$ to a terminal, so $l\le n$. Let $y_0=x_v$, and for $i\ge1$ let
$y_i$ be the load on the edge(s) from $u_{i-1}$ to $u_i$. Then $y_i\le x_{u_i}$ for $i<l$ and
$y_l\le1$. Since $y_i\ge\phi(x_{u_{i-1}})$ and $\phi$ is nondecreasing, $y_i\ge\phi(y_{i-1})$.
Induction gives $y_i\ge(K+1)/2^i-1$, because $(K+1)/2^{i+1}-1<K/2$. So
$y_l\ge(2^{n+1}+1)/2^n-1=1+2^{-n}>1$, which contradicts the bound on terminal inflow. Hence $x_v\le K$
everywhere, S1 gives $A_v=h_0(x_v)$ and $B_v=h_1(x_v)$, and the load equations are exactly the
fixed-point equations.

\emph{(d)} Property 3 then follows from Lemma~\ref{lem:decoder}.
\end{proof}

\section{The switch gadget}\label{sec:ar-gadget}

\subsection{The counting network}

\paragraph{Idea.} In the counting network $G_k$, the $y$-th unit of flow sent along a cheapest path
costs the $y$-th entry of a list $D_k$. The list grows by mirroring: $D_k$ is $D_{k-1}$ followed by its
reflection about $3^k$. All gaps are at least $6$, so there are no ties when $0<\theta<6$ (Lemma~\ref{lem:alt}).

$G_0$ is a single arc $s_0\to t_0$ with capacity $1$ and cost $0$. $G_k$ adds nodes $s_k,t_k$ and four
arcs of capacity $2^{k-1}$: $s_k\to s_{k-1}$ and $t_{k-1}\to t_k$ with cost $0$, and $s_k\to t_{k-1}$
and $s_{k-1}\to t_k$ with cost $3^k$ (Figure~\ref{fig:Gk}). This is the counting gadget of Disser and
Skutella~\cite[\S3.1]{ds15}, whose expensive arcs cost $2^{k-1}-\tfrac12$; the cost $3^k$ widens the gaps
between increments.

\begin{figure}[t]
\centering
\begin{tikzpicture}[
  nd/.style={circle,draw,inner sep=1.5pt,minimum size=6.5mm,font=\small},
  ar/.style={-{Stealth[length=2.2mm]},thick},
  ex/.style={-{Stealth[length=2.2mm]},thick,red!70!black},
  lb/.style={font=\scriptsize,fill=white,inner sep=1pt}]
  \node[nd] (s0) at (0,0) {$s_0$};
  \node[nd] (t0) at (5,0) {$t_0$};
  \node[nd] (s1) at (0,2) {$s_1$};
  \node[nd] (t1) at (5,2) {$t_1$};
  \node[nd] (s2) at (0,4.4) {$s_2$};
  \node[nd] (t2) at (5,4.4) {$t_2$};
  \draw[ar] (s0) -- node[lb,below=1pt]{$1,\ 0$} (t0);
  \draw[ar] (s1) -- node[lb,left=1pt]{$1,\ 0$} (s0);
  \draw[ar] (t0) -- node[lb,right=1pt]{$1,\ 0$} (t1);
  \draw[ex] (s1) -- node[lb,pos=0.22]{$1,\ 3$} (t0);
  \draw[ex] (s0) -- node[lb,pos=0.22]{$1,\ 3$} (t1);
  \draw[ar] (s2) -- node[lb,left=1pt]{$2,\ 0$} (s1);
  \draw[ar] (t1) -- node[lb,right=1pt]{$2,\ 0$} (t2);
  \draw[ex] (s2) -- node[lb,pos=0.22]{$2,\ 9$} (t1);
  \draw[ex] (s1) -- node[lb,pos=0.22]{$2,\ 9$} (t2);
  \node[font=\small,align=left,anchor=west] at (6.2,2.2) {labels: capacity, cost\\[2pt]
    \textcolor{red!70!black}{red}: expensive arcs\\[2pt]
    $D_1=(0,6)$\\ $D_2=(0,6,12,18)$};
\end{tikzpicture}
\caption{The counting network $G_2$ (two levels). Level $k$ adds the arcs $s_k\to s_{k-1}$ and
$t_{k-1}\to t_k$ of cost $0$, and the crossing arcs $s_k\to t_{k-1}$ and $s_{k-1}\to t_k$ of cost $3^k$;
all four have capacity $2^{k-1}$.}
\label{fig:Gk}
\end{figure}

\begin{lemma}[Increments]\label{lem:incr}
\statusrefereed
Let $F_k(y)$ be the minimum cost of an $s_k\to t_k$ flow of value $y\in[0,2^k]$. Then $F_k$ is convex and
piecewise linear with integer breakpoints. Its list of unit increments satisfies $D_0=(0)$ and
\[
D_k=\bigl(D_{k-1},\;2\cdot3^k-\mathrm{reverse}(D_{k-1})\bigr).
\]
For $k\ge1$ the list is strictly increasing, starts at $0$, ends at $2\cdot3^k$, and has minimum gap
exactly $6$. The maximum flow is $2^k$.
\end{lemma}

\begin{proof}
Induction on $k$. Put $c=2^{k-1}$ and $w=3^k$, and let $\alpha,\beta,\gamma,\eta$ be the flows on
$s_k\to s_{k-1}$, $t_{k-1}\to t_k$, $s_k\to t_{k-1}$ and $s_{k-1}\to t_k$. The inner $G_{k-1}$ is acyclic
and meets the outer arcs only at $s_{k-1}$ and $t_{k-1}$, so it carries an $s_{k-1}\to t_{k-1}$ flow of
some value $y$. Conservation gives $\alpha+\gamma=x$, $\alpha=y+\eta$, $\beta=y+\gamma$ and
$\beta+\eta=x$, so $\gamma+\eta=x-y$. The capacity box is feasible if and only if
$0\le y\le\min(x,2c-x)$. So $F_k(x)=\min_y[wx-wy+F_{k-1}(y)]$ over that range. Every increment of
$F_{k-1}$ is at most $\max D_{k-1}$, which is $0$ for $k=1$ and $2\cdot3^{k-1}$ for $k\ge2$, and both are
below $3^k$. So the bracket strictly decreases in $y$, and $y=\min(x,2c-x)$. This gives $F_k=F_{k-1}$ on
$[0,c]$ and $F_k(x)=2w(x-c)+F_{k-1}(2c-x)$ on $[c,2c]$, which yields the recursion for $D_k$. The gap at
the junction is $2w-2\max D_{k-1}$. It equals $6$ for $k=1$, because $D_0=(0)$, and $2\cdot3^{k-1}\ge6$
for $k\ge2$. All other gaps are gaps of $D_{k-1}$. The value $x=2c$ is feasible.
\end{proof}

For example, $D_1=(0,6)$, $D_2=(0,6,12,18)$ and $D_3=(0,6,12,18,36,42,48,54)$.

\subsection{Two copies and the alternation}

Let $N_m$ consist of two copies of $G_m$ that share their source $s$. Copy A ends at sink $a$, and copy
B ends at sink $b$. Each unit entering $b$ costs an extra $\theta$. The copies share only $s$. Each sink
has in-capacity $2^m$, and $s$ has out-capacity $K=2^{m+1}$ (Figure~\ref{fig:Nm}).

\begin{figure}[t]
\centering
\begin{tikzpicture}[
  nd/.style={circle,draw,inner sep=1.5pt,minimum size=7mm,font=\small},
  bx/.style={draw,rounded corners,minimum width=3.4cm,minimum height=1.1cm,align=center,font=\small},
  ar/.style={-{Stealth[length=2.2mm]},thick}]
  \node[nd] (s) at (0,0) {$s$};
  \node[bx] (A) at (3.6,1.3) {copy A of $G_m$};
  \node[bx] (B) at (3.6,-1.3) {copy B of $G_m$};
  \node[nd] (a) at (9.6,1.3) {$a$};
  \node[nd] (b) at (9.6,-1.3) {$b$};
  \draw[thick] (s) -- (A.west);
  \draw[thick] (s) -- (B.west);
  \draw[ar] (A.east) -- node[above,font=\scriptsize]{in-capacity $2^m$} (a);
  \draw[ar] (B.east) -- node[above,font=\scriptsize]{in-capacity $2^m$}
     node[below,font=\scriptsize]{extra cost $\theta$ per unit} (b);
  \node[font=\scriptsize,align=center,anchor=east] at (-0.4,0) {out-capacity\\$K=2^{m+1}$};
\end{tikzpicture}
\caption{The splitter $N_m$: two copies of the counting network share their source $s$ (the node $s_m$
of both copies); copy A's sink $t_m$ is $a$ and copy B's is $b$. Copy B is offset by $\theta$ per unit.}
\label{fig:Nm}
\end{figure}

\paragraph{Idea.} Copy A's $y$-th unit costs $\delta_y$, copy B's costs $\delta_y+\theta$. With
$0<\theta<6$ and all gaps of the increment list at least $6$, the merged order of unit costs is
$\delta_1<\delta_1+\theta<\delta_2<\delta_2+\theta<\dots$, so the cheapest flow sends units to
$a,b,a,b,\dots$. Table~\ref{tab:alt} shows this for $m=2$ and $\theta=3$.

\begin{table}[h]
\centering
\small
\begin{tabular}{@{}lcccccccc@{}}
\toprule
unit number & 1 & 2 & 3 & 4 & 5 & 6 & 7 & 8\\
\midrule
cheapest remaining unit & $0$ & $3$ & $6$ & $9$ & $12$ & $15$ & $18$ & $21$\\
goes to & $a$ & $b$ & $a$ & $b$ & $a$ & $b$ & $a$ & $b$\\
$\inflow(a)=h_0(x)$ after $x$ units & 1 & 1 & 2 & 2 & 3 & 3 & 4 & 4\\
$\inflow(b)=h_1(x)$ after $x$ units & 0 & 1 & 1 & 2 & 2 & 3 & 3 & 4\\
\bottomrule
\end{tabular}
\caption{Alternation in $N_2$ with $\theta=3$. Copy A's increments are $D_2=(0,6,12,18)$; copy B's
are $D_2+3=(3,9,15,21)$.}
\label{tab:alt}
\end{table}

\begin{lemma}[Alternation]\label{lem:alt}
\statusrefereed
Let $0<\theta<6$ and $x\in[0,K]$. Among the flows of value $x$ from $s$ to $\{a,b\}$, every flow
minimizing $\sum_e\kappa_ef_e+\theta\cdot\inflow(b)$ has $\inflow(a)=h_0(x)$ and $\inflow(b)=h_1(x)$.
\end{lemma}

This is Disser and Skutella's alternation argument~\cite[\S3.2, proof of Lemma 2]{ds15}. Their network
joins two copies of the counting gadget at a common source and a common sink, and they show that
successive shortest paths alternate between the copies. The version here differs in four ways: the sinks
are separate, there is no cross arc, a constant offset puts copy A first in each pair, and the gaps are
wide enough that the order survives every pair of port prices in the box.

\begin{proof}
A flow splits as $A+B=x$ with $0\le A,B\le2^m$, and its cost is at least
$F_m(A)+F_m(B)+\theta B$, attained by cheapest flows for both parts.
Each copy offers $2^m$ unit intervals. If $\delta_j$ is the $j$-th
marginal cost of copy A, the corresponding cost in B is $\delta_j+\theta$.
The gap bound and $0<\theta<6$ strictly order the combined list as
$\delta_1,\delta_1+\theta,\delta_2,\delta_2+\theta,\ldots$.
Filling its cheapest $x$ units respects the interval order in each
copy. Any different split uses some dearer interval while a cheaper
one still has room, and transferring a positive amount lowers the cost.
The unique split therefore gives $A=h_0(x)$ and $B=h_1(x)$, including
$x=0$ and $x=K$.
\end{proof}

At $\theta=0$ or $\theta=6$ ties appear. For example, at $\theta=0$ every split is optimal for
$x\in(0,1]$, since both copies have $\delta_1=0$. This is why condition C2 below is strict.

\subsection{Market embedding of the gadget}\label{sec:ar-gadget-market}

\paragraph{Idea.} Turn the network into a market: arcs become goods, nodes become buyers who
``buy'' arcs, and arc costs become payoff differences. With unit slopes this market is a linear program
(Lemma~\ref{lem:lp}), and its optimal flows are min-cost flows of $N_m$ with $\theta=\varepsilon+P_b-P_a$.

Let the arcs of $N_m$ be $e=(u,w)$, with capacity $c_e$ and cost $\kappa_e$.
\begin{itemize}[nosep]
\item Each arc becomes a private good $g_e$ with capacity $c_e$.
\item The tail $u$ holds an option on $g_e$ with payoff $M-\kappa_e-P_{g_e}$. Its purchase $f_e$ is the
arc flow.
\item The head $w$ holds an option on $g_e$ with payoff $M-P_{g_e}$ if $w$ is internal, and
$M+H-P_{g_e}$ if $w\in\{a,b\}$. Its purchase $\sigma_e$ is slack.
\item An internal node is a buyer whose mass is the total capacity of its in-arcs.
\item Sink $a$ (respectively $b$) is a buyer of mass $2^m$. It also has a port option with payoff
$M-P_a$ (respectively $M-\varepsilon-P_b$).
\item The source $s$ is a buyer of mass $c_p$. Its options are $p$, with payoff $M_p-\lambda P_p$, and
its out-arc goods.
\item The input good $p$ has capacity $c_p$.
\end{itemize}

We use $\varepsilon=3$, $H=1$, $\bar P=1$, $M=4(\sum_ec_e\kappa_e+\varepsilon+H+10)$, $M_p=M+1$ and
$\lambda=M_p/\bar P$. Standalone, $c_p=2K$. In the composition, $c_{p_v}=\max(2K,\ \indeg(v)\cdot2^m+2)$,
where $\indeg$ counts labelled in-edges. Each port buyer has mass $2^m=K/2$, so this choice meets S3.
The terminals use $C_T=\indeg(T)\cdot2^m+2$, a terminal buyer of mass $C_T+1$ with payoff $1/2-P_T$ (so
$r_T=1/2$), and a dispatch buyer with payoff $2-P_{p_o}$. The only slope other than $1$ is $\lambda$, on
the source's option on $p$, so every slope is at least $1$, as Theorem~\ref{thm:mk-embed} requires. S4
holds because $s$ holds $p$, sink $a$ holds $a$ and sink $b$ holds $b$.

The proof uses the following conditions:
\begin{description}[leftmargin=2.5em,itemsep=0pt]
\item[C1:] $M_p/\lambda\le\bar P$;
\item[C2:] $\bar P<\varepsilon<6-\bar P$, with strict inequalities;
\item[C3:] $H>0$;
\item[C4:] $M>\varepsilon+\bar P+\Upsilon$, where $\Upsilon=\sum_e\kappa_e+4H=2\cdot3^{m+1}-6+4H$;
\item[C5:] $M_p>M-H$;
\item[C6:] $c_p>K$, and buyer $s$ has mass $c_p$.
\end{description}
The parameters above satisfy all six.

\begin{theorem}[The gadget meets the specification]\label{thm:gadget}
\statusrefereed
Under C1--C6, for every $m\ge1$, the gadget meets S1 and S2 with $K=2^{m+1}$ and box $[0,\bar P]^2$.
\end{theorem}

The proof has three steps: the gadget market is an LP (Lemma~\ref{lem:lp}); every private good is full
(Lemma~\ref{lem:full}); and every optimum of the LP is a flow of value $\min(x,K)$
(Lemma~\ref{lem:flow}). Then Lemma~\ref{lem:alt} finishes.

\begin{lemma}[LP form]\label{lem:lp}
Fix $x\in[0,c_p]$ and port prices $P_a,P_b\ge0$, and put $\pi=\lambda P_p$. Consider the welfare LP (W).
It maximizes $\sum\tilde a_{ij}f_{ij}$ subject to the row bounds (mass) and column bounds (capacity of
private goods, and $c_p-x$ for $p$). Here $\tilde a=a_{ij}$ on private goods, $\tilde a=M_p$ on $(s,p)$,
and $\tilde a=a_{it}-P_t$ on port options, which have no column. Then gadget equilibria are exactly the
optimal primal--dual pairs of (W), with dual prices $(P_{\text{private}},\pi)$.
\end{lemma}

\begin{proof}
With unit slopes on private goods, conditions (a)--(d) are primal feasibility, dual feasibility and
complementary slackness. In particular $v_s\ge M_p-\lambda P_p$ is the same as $v_s+\pi\ge M_p$.
\end{proof}

The LP form has three consequences. Equilibria exist at all port prices $\ge0$. Every optimal flow pairs
with every optimal dual. And if $x<c_p$ and $P_p>0$, then $s$ buys $p$, which gives
$P_p\le M_p/\lambda\le\bar P$.

\begin{lemma}[Private goods are full]\label{lem:full}
Assume C3, $\varepsilon\ge0$ and $M>0$, and take any $x$ and any $P_a,P_b\ge0$. In every equilibrium, $f_e+\sigma_e=c_e$ for every
arc $e$.
\end{lemma}

\begin{proof}
Every internal buyer other than the source $s$ has payoffs at most $M$; the walk below visits only heads of arcs, and $s$ is never one. A sink's payoffs are at most $M+H$, and its port option pays
less than $M+H$.

Suppose first that an in-arc of sink $t$ has price $0$. Then $v_t=M+H$, so $t$ is full and buys only
zero-price slack. Its mass equals its in-capacity, so $\sigma_e=c_e$ and $f_e=0$ on all its in-arcs; call
this $(\ast)$.

Now suppose $e=(u,w)$ with $w$ internal is not full. Then $P_{g_e}=0$, so $v_w=M$ and $w$ is full. Its
slack is less than its mass, so it buys flow on an out-arc with $\kappa=0$ and price $0$. The head of that
arc is again at value $M$, and the argument repeats. $N_m$ is acyclic, so the walk reaches a sink
through a zero-price arc with positive flow, which contradicts $(\ast)$.

An arc into a sink that is not full has price $0$, so it is full by $(\ast)$, a contradiction.
\end{proof}

\paragraph{Flow notation.} Write $\varepsilon_a=0$ and $\varepsilon_b=\varepsilon$ for the port offsets of the two
sinks, so sink $t$'s port option pays $M-\varepsilon_t-P_t$. For a flow $f$ on $N_m$, $F_{\mathrm{out}}$ is
the total flow on the out-arcs of $s$, the \emph{excess} of an internal node is its inflow minus its
outflow, and the \emph{residual network} has a forward arc $e$ wherever $f_e<c_e$ and a backward copy of
$e$ wherever $f_e>0$. Pushing $\delta$ along a residual path raises the flow on its forward arcs and lowers
it on its backward arcs, so its cost change is at most $\delta$ times the sum of $\kappa'_e$ over its
forward arcs.

\begin{lemma}[Flow structure]\label{lem:flow}
Assume C3--C6, $\varepsilon\ge0$, $x\in[0,c_p)$ and $(P_a,P_b)\in[0,\bar P]^2$ (the construction uses $\varepsilon=3$). Substitute $\sigma_e=c_e-f_e$, which is
valid by Lemma~\ref{lem:full}, and write $z_t$ for sink $t$'s port purchase and $y$ for $s$'s purchase
of $p$. Up to a constant, (W) becomes the program (R):
\[
\text{maximize } M_py-\sum_e\kappa'_ef_e+\sum_tg_tz_t,
\]
where $\kappa'_e=\kappa_e+H\cdot[\text{head is a sink}]$ and $g_t=M-\varepsilon_t-P_t$, subject to
$y\le c_p-x$, $y+F_{\mathrm{out}}\le c_p$, outflow $\le$ inflow at internal nodes, $z_t\le\inflow(t)$
and $0\le f\le c$. Every equilibrium flow is optimal for (R), and every optimum of (R) has $y=c_p-x$,
conservation at every internal node, $z_t=\inflow(t)$ and $F_{\mathrm{out}}=\min(x,K)$.
\end{lemma}

\begin{proof}
By C4, $g_t\ge M-\varepsilon-\bar P>\Upsilon\ge\sum_e\kappa'_e$.
\begin{enumerate}[nosep]
\item If $z_t<\inflow(t)$, raising $z_t$ gains $g_t>0$.
\item Suppose $y<c_p-x$. If $s$ has spare mass, raising $y$ gains $M_p$. Otherwise
$F_{\mathrm{out}}>x\ge0$. Take a positive-flow path from $s$ to a sink or to an internal node with excess,
and move $\delta$ from it to $y$, lowering $z_t$ if the path ends at sink $t$. The gain is at least
$\delta(M_p-\max(0,M-H))>0$ by C5. So $y=c_p-x$, and therefore $F_{\mathrm{out}}\le x$.
\item Suppose an internal node $w$ has excess, and let $R$ be its residual-reachable set. If $R$ contains
a sink $t$, push $\delta$ along a simple residual path to it and raise $z_t$ by $\delta$. The gain is at
least $\delta(g_t-\Upsilon)>0$, because a simple path costs at most the sum of $\kappa'$ over its forward
arcs. Otherwise, walking back from $w$ along positive arcs reaches $s$, so $s\in R$. Every arc leaving
$R$ is saturated and every arc entering $R$ is empty. So
$\mathrm{cap}(\delta^+(R))=F_{\mathrm{out}}-\sum_{u\in R}\mathrm{excess}(u)<F_{\mathrm{out}}\le K$. But
$R$ separates $s$ from both sinks, so its capacity is at least the maximum flow $K$
(Lemma~\ref{lem:incr}). Contradiction.
\item If $F_{\mathrm{out}}<\min(x,K)$, then $s$ has spare mass. Either a residual path from $s$ to a
sink exists, and pushing along it gains $g_t-\Upsilon>0$, or there is a cut of capacity
$F_{\mathrm{out}}<K$. Both are impossible.\qedhere
\end{enumerate}
\end{proof}

\begin{proof}[Proof of Theorem~\ref{thm:gadget}]
By Lemmas~\ref{lem:lp}--\ref{lem:flow}, an equilibrium flow is an $s\to\{a,b\}$ flow of value
$\min(x,K)$, with $A=\inflow(a)$ and $B=\inflow(b)$. On such flows the objective of (R) equals a
constant minus $[\sum_e\kappa_ef_e+\theta B]$, where $\theta=\varepsilon+P_b-P_a$.
\begin{itemize}[nosep]
\item S1: for $x\le K$, C2 gives $\theta\in[\varepsilon-\bar P,\varepsilon+\bar P]\subset(0,6)$, and
Lemma~\ref{lem:alt} gives $A=h_0(x)$ and $B=h_1(x)$. The only $p$-labelled option is $(s,p)$, whose
ratio is $M_p/\lambda\le\bar P$ by C1.
\item S2: for $x\in[K,c_p)$ we have $A+B=K\le x$ and $A,B\le2^m$, so $A=B=K/2$.\qedhere
\end{itemize}
\end{proof}

\paragraph{Sharpness.} Exact computations confirm that C2 cannot be weakened to closed inequalities
and that C3 cannot be dropped. With $\varepsilon\in\{1,5\}$
the box reaches the ties $\theta\in\{0,6\}$. With $H=0$, sink $a$ buys its port with no flow behind it at
$x=0$. The strict inequality in C6 is convenient for the composition, where S3 ensures $x<c_p$ and hence
caps $P_p$; at $c_p=K$, $P_p$ can be unbounded at $x=K$, although the stated S1/S2 output properties still
hold. (The exact checker also tests the input-price bound $P_p\le\bar P$, which is stronger than S1; that
check, not S1, fails at $c_p=K$.) C4 is far from tight, and C5 was not probed.

\paragraph{Where the non-unit slope matters.} If every slope is $1$, the (M1) conditions are the
optimality conditions of one transportation LP. The same holds after rescaling when the slopes factor as
$b_{ij}=\beta_i\gamma_j$. So with such slopes a reduction of this kind would put ARRIVAL in P. Inside one gadget at
fixed port prices, $\lambda$ only rescales the price of $p$ (Lemma~\ref{lem:lp}), and the gadget is an LP.
Across gadgets, the good $p_v$ is bought at slope $\lambda$ by $v$'s source and at slope $1$ by the port
buyers of $v$'s predecessors. All buyers of one gadget are linked through slope-$1$ private goods, so a
factorization would need $\beta_u=\beta_v/\lambda$ along every edge $u\to v$ of the switch graph. That
is impossible around a cycle. We know no reduction of these markets to an LP or to a concave generalized
flow problem~\cite{vegh14}. Such a reduction, combined with this chapter, would put ARRIVAL in P.

\section{Sizes, and the proof of the theorem}\label{sec:ar-size}

\begin{itemize}[nosep]
\item \emph{Gadget of depth $m$.} The network $N_m$ has $8m+2$ arcs, and the gadget has $4m+3$ buyers.
Counting buyers, private goods and $p$ gives $12m+6=12\log_2K-6$ buyers and goods. The numbers are capacities
$2^{k-1}$, costs $3^k$ and $M=\Theta(6^m)$, so there are $O(m)$ numbers of $O(m)$ bits each. Measured
encodings take $282$, $564$, $978$, $1428$, $2042$, $2660$ and $3474$ bits for $K=4,\dots,256$. The
exponentially large quantities sit in binary-encoded capacities and intercepts, so algorithms that are
pseudo-polynomial in these numbers are exponential on these instances.
\item \emph{Composed market ($m=n$, $K=2^{n+1}$).} The market has $4n^2+3n+3$ buyers and $8n^2+3n+2$
goods. Every number has $O(n)$ bits. The largest are $M$ and $\lambda=M+1$, of order $6^n$, and
$c_{p_v}=O(n2^n)$. For $n=2$, $M=392$ and $\lambda=393$.
\item \emph{Game.} The game has at most $12n^2+6n+6$ types and at most $12n^2+6n+7$ targets (fewer if preprocessing
deletes vertices). The constants $\hat D$, $W$, $A$ and $S$ of~\eqref{eq:mk-consts}, and every payoff and
prior, have $O(n)$ bits. Listing payoffs for every (type, target) pair takes $O(n^4)$ numbers, or
$O(n^5)$ bits. For $n=1$ there are $24$ types and $25$ targets, and for $n=2$ there are $66$ types and
$67$ targets, with $A=394$ and $S=25611$. The constant $W$ depends on the in-degrees: for the checked
two-vertex instance with $(\succz(0),\succo(0))=(0,1)$ and $(\succz(1),\succo(1))=(\yes,\no)$ it is
$W=400$. For the chain instance with $n=7$, $\succz(v)=v+1$ (with $\succz(6)=\yes$), $\succo(v)=0$ and
origin $0$, the explicit game takes about $2.6$ Mbit, counting numerator and denominator bits of every
prior and every listed payoff.
Both $\Red$ and $\Dec$ run in polynomial time.
\end{itemize}

The market part of the construction is worth recording on its own, since Chapter~\ref{ch:upper} uses it
without the game.

\begin{theorem}[ARRIVAL to (M1)]\label{thm:arrival-m1}
\statusrefereed
There is a polynomial-time map from ARRIVAL instances with $n$ switch vertices to (M1) markets with the
following properties. After a preprocessing step that deletes vertices unable to reach a terminal (which
preserves the answer), the market is composed of one copy, for each vertex $v$, of a switch gadget built
from two copies of the Disser--Skutella counting network, with an input good $p_v$, plus a dispatch buyer
of mass $1$ on $p_o$, and for each terminal $t$ a terminal good $T_t$ and a terminal buyer. All masses
and capacities are positive and all slopes are at least $1$; the market has at most $4n^2+3n+3$
buyers, at most $8n^2+3n+2$ goods, and numbers of $O(n)$ bits. The construction satisfies the gadget specification
(S1--S4) with $K=2^{n+1}$, and each terminal capacity $C_T$ exceeds the total mass of the non-terminal
buyers holding an option on $T_t$. Under these conditions, in \emph{every} equilibrium of the market the
load placed on $T_t$ by non-terminal buyers is $1$ for the terminal $t$ that the run of the preprocessed instance reaches and $0$ for
the other terminal.
\end{theorem}

\begin{proof}
Preprocess (Section~\ref{sec:pre-arrival}) and compose the gadget of
Section~\ref{sec:ar-gadget-market} with $m=n$. By Theorem~\ref{thm:gadget} it meets S1 and S2; S3 and
the terminal condition hold by the choice of $c_{p_v}$ and $C_T$ above; S4 holds by the structure of the
gadget ($s$ holds $p$, sink $a$ holds $a$, sink $b$ holds $b$). The dispatch buyer has ratio $2>\bar P=1$,
and the terminal buyers have ratio $r_T=1/2\in(0,\bar P)$, as the composition requires. The conclusion is
Theorem~\ref{thm:compose}(3), and the sizes are those just listed.
\end{proof}

\begin{proof}[Proof of Theorem~\ref{thm:A}]
$\Red$ builds the market of Theorem~\ref{thm:arrival-m1} and then the game $\Gamma_1$ of
Theorem~\ref{thm:mk-embed} (with $\omega=1$, so $W=4\hat D$); this is legal because every slope is $1$
or $\lambda\ge1$. Given any BNE, $\Dec$ decodes it to a market equilibrium by Theorem~\ref{thm:mk-embed}
and answers \yes{} if and only if the non-terminal buyers' load on $T_{\yes}$ is $1$, which is correct
by Theorem~\ref{thm:arrival-m1}. The game has one resource and defender payoffs $0$ and $1$, and its size
is as listed above.
\end{proof}

\begin{lemma}[No mass creation at arbitrary port prices]\label{rem:ar-smallterm}
For the two-copy gadget of Section~\ref{sec:ar-gadget-market}, at any
nonnegative port prices and any $0\le x<c_p$, its port purchases satisfy
$A+B\le x$. Consequently any terminal capacity $C_T>1$ suffices in
Theorem~\ref{thm:compose} for this construction, with the same terminal
buyer and the same switch input capacities.
The numerical choices in Section~\ref{sec:ar-gadget-market} give $M_p=M+1$.
\end{lemma}
\noindent{\small\itshape AI self-review, November 26, 2026 (UTC).
Checked for arbitrary nonnegative port prices, including a negative net payoff at a port.}\par
\begin{proof}
Use the residual assignment LP of Lemma~\ref{lem:lp}. Lemma~\ref{lem:full}
already applies at all nonnegative port prices. Substituting the head
purchase $c_e-f_e$ therefore gives the objective and feasible inequalities
of Lemma~\ref{lem:flow}, although the rewards $g_t=M-\varepsilon_t-P_t$
can now be arbitrarily negative. We need only the upper bound $g_t\le M$.
Let $y$ be the source's input purchase. If $y<c_p-x$ and the source has
spare mass, increasing $y$ strictly improves the objective. Otherwise the
source is full and has a positive outgoing flow. Follow positive arcs to
the first internal vertex with surplus inflow or to a sink. All intermediate
vertices can pass positive flow onward, and acyclicity makes the path finite.
Remove a sufficiently small amount $\eta>0$ along it and add $\eta$ to $y$.
If the path ends at a sink, decrease its port purchase only as far as needed
to maintain $z_t\le\inflow(t)$. This costs at most $\eta\max(0,g_t)\le\eta M$,
and all saved arc fees are nonnegative. The entry gains $\eta M_p>\eta M$.
Optimality forces $y=c_p-x$. Thus the source outflow is at most $x$.
Each internal outflow is at most its inflow and each sink purchase is at
most its inflow. Summing yields $A+B\le x$.

For the composition, $x_v<c_{p_v}$ follows from the switch input capacities
regardless of any terminal price. A positive input price needs a source
purchase, so it is at most $1$, and the single dispatch buyer spends its unit.
Apply the inequality just proved at every switch and cancel all internal
port purchases. The sum of external receipts at the terminals is at most
$1$. If $P_T>1/2$, its dedicated buyer leaves and a terminal of capacity
$C_T>1$ cannot fill. If $P_T<1/2$, that buyer alone demands more than $C_T$.
Hence $P_T=1/2$, which puts every port in the required price box and
completes the original proof. The earlier theorem's stated size bounds
also follow using its original larger capacities.
\end{proof}
